
\documentclass[11pt]{article}
\usepackage[margin=1in]{geometry}
\usepackage[T1]{fontenc}
\usepackage[utf8]{inputenc}
\usepackage{lmodern}
\usepackage{amsmath}
\usepackage{amssymb}
\usepackage{graphicx}
\usepackage{booktabs}
\usepackage{float}
\usepackage{hyperref}
\hypersetup{hidelinks}
\setlength{\parskip}{0.5em}
\setlength{\parindent}{0pt}
\begin{document}
\sloppy
\section{Step 1: Harmonization and Descriptive Analysis of PsychoPy Demographic Files}
\label{sec:step1_demographics}
% Analyst note added 2026-03-20: PID006 and PID037 are globally excluded from downstream participant-level analyses. PID006 was excluded after Step 3 documented session 2025-11-12_002 as an implausibly short 10.420224 second recording with no paired PsychoPy log. PID037 was excluded because a behavioral log exists but no fNIRS session folder was found. Step 1 outputs are regenerated without these participants, and future steps should honor materials/analysis_participant_exclusions.json.

\subsection{Objective}
The objective of Step~1 is to create a clean, traceable, and fully documented demographic dataset
from the PsychoPy outputs before any question-level or fNIRS-level analysis is performed.
This step focuses only on the demographic and socioeconomic information collected before the task,
together with the questionnaire metadata needed to interpret those responses correctly.

The present study investigates abstract versus concrete exam questions using fNIRS in an educational
setting, and the project documentation states that participant information is pseudonymized and used
in aggregated form. The demographics component therefore needs to be clean, consistent, and easy
to join later with behavioral and fNIRS data. \cite{step1_project_doc}

\subsection{Current Raw Data Structure}
Each participant appears only once in the dataset.
Each participant has a dedicated folder named according to the pattern:
\begin{verbatim}
PID001/
PID002/
...
\end{verbatim}

Inside each participant folder, at minimum, the following files are expected:
\begin{verbatim}
enem_blocks_PID01_20251106_124958.csv
socio_PID01_20251106_124958.json
\end{verbatim}

The \texttt{enem\_blocks} file contains event-level PsychoPy information, including the start of the
questionnaire block and questionnaire item markers.
The \texttt{socio} JSON file contains the actual demographic answers keyed by variable name.

In addition, the file \texttt{demographics\_questions.json} defines the questionnaire schema, including
question identifiers, prompts, variable types, required flags, and conditional display rules.
This file should be treated as the primary questionnaire specification for Step~1. \cite{step1_demo_schema}

\subsection{Principles for Step~1}
Step~1 should follow six basic principles:
\begin{enumerate}
    \item Keep the workflow simple and conservative.
    \item Never overwrite or modify any raw file.
    \item Preserve both raw values and harmonized values.
    \item Treat \texttt{demographics\_questions.json} as the reference schema for current sessions.
    \item Flag ambiguous legacy values rather than forcing recodes.
    \item Produce modular outputs: CSV tables, plots, and \LaTeX{} snippets generated automatically.
\end{enumerate}

\subsection{Scope}
This step includes:
\begin{itemize}
    \item inventory of demographic-related raw files;
    \item extraction and validation of participant identifiers and timestamps;
    \item harmonization of demographic variables;
    \item grouped descriptive analysis of each demographic question;
    \item export of clean CSV files;
    \item generation of plots for each variable;
    \item generation of one \LaTeX{} file per variable and one master \LaTeX{} report.
\end{itemize}

This step does not include:
\begin{itemize}
    \item analysis of question performance;
    \item reconstruction of missing fNIRS triggers;
    \item preprocessing or statistical analysis of the fNIRS signals.
\end{itemize}

\subsection{Authoritative Variable List}
The current questionnaire schema contains the following variables:
\begin{verbatim}
age
country_of_birth
nationality
years_in_luxembourg
gender
first_language
other_languages
affiliation
department_unit
programme_name
study_level
year_in_programme
employee_role
highest_degree
degree_field
country_secondary_ed
primary_school_langs
parent_a_edu
parent_b_edu
first_gen_he
living_arrangement
\end{verbatim}

Some of these variables are conditional.
For example, \texttt{programme\_name}, \texttt{study\_level}, and \texttt{year\_in\_programme}
are displayed only if affiliation matches a student pattern, while \texttt{employee\_role} is shown
only for employee/staff/faculty/postdoc patterns. These conditions are explicitly defined in the
questionnaire schema and must be respected during analysis. \cite{step1_demo_schema}

\subsection{Question Groups}
For organization and reporting, the variables should be analyzed in the following groups:

\begin{enumerate}
    \item \textbf{Core identification and mobility variables}
    \begin{itemize}
        \item \texttt{age}
        \item \texttt{country\_of\_birth}
        \item \texttt{nationality}
        \item \texttt{years\_in\_luxembourg}
    \end{itemize}

    \item \textbf{Gender and language background}
    \begin{itemize}
        \item \texttt{gender}
        \item \texttt{first\_language}
        \item \texttt{other\_languages}
        \item \texttt{primary\_school\_langs}
    \end{itemize}

    \item \textbf{Current academic or professional status}
    \begin{itemize}
        \item \texttt{affiliation}
        \item \texttt{department\_unit}
        \item \texttt{programme\_name}
        \item \texttt{study\_level}
        \item \texttt{year\_in\_programme}
        \item \texttt{employee\_role}
    \end{itemize}

    \item \textbf{Educational background}
    \begin{itemize}
        \item \texttt{highest\_degree}
        \item \texttt{degree\_field}
        \item \texttt{country\_secondary\_ed}
    \end{itemize}

    \item \textbf{Family educational background and living conditions}
    \begin{itemize}
        \item \texttt{parent\_a\_edu}
        \item \texttt{parent\_b\_edu}
        \item \texttt{first\_gen\_he}
        \item \texttt{living\_arrangement}
    \end{itemize}
\end{enumerate}

\subsection{Expected Inputs}
The minimum raw inputs for Step~1 are:
\begin{itemize}
    \item one participant folder per participant, e.g. \texttt{PID001/};
    \item one PsychoPy event file per participant:
    \texttt{enem\_blocks\_PIDxx\_YYYYMMDD\_HHMMSS.csv};
    \item one demographic response file per participant:
    \texttt{socio\_PIDxx\_YYYYMMDD\_HHMMSS.json};
    \item the questionnaire schema file:
    \texttt{demographics\_questions.json}.
\end{itemize}

\subsection{Expected Outputs}
The following outputs should be generated in Step~1:

\subsubsection*{Structured tables}
\begin{verbatim}
01_file_manifest.csv
02_participant_registry.csv
03_demographics_schema.csv
04_socio_raw_wide.csv
05_socio_raw_long.csv
06_socio_harmonization_dictionary.csv
07_socio_clean_wide.csv
08_question_summary.csv
09_question_level_counts.csv
10_question_missingness.csv
\end{verbatim}

\subsubsection*{Plots}
\begin{verbatim}
figures/step1/age_histogram.png
figures/step1/years_in_luxembourg_histogram.png
figures/step1/gender_barplot.png
figures/step1/affiliation_barplot.png
...
\end{verbatim}

\subsubsection*{\LaTeX{} outputs}
\begin{verbatim}
reports/step1/questions/question_age.tex
reports/step1/questions/question_gender.tex
reports/step1/questions/question_affiliation.tex
...
reports/step1/step1_demographics_report.tex
\end{verbatim}

\subsection{Folder Structure for Step~1}
A simple and organized structure is recommended:
\begin{verbatim}
project/
+-- data_raw/
|   +-- PID001/
|   +-- PID002/
|   +-- ...
+-- data_intermediate/
|   +-- step1/
+-- data_clean/
|   +-- step1/
+-- figures/
|   +-- step1/
+-- reports/
|   +-- step1/
|       +-- questions/
+-- scripts/
\end{verbatim}

\subsection{Detailed Workflow}

\subsubsection{Step~1.1: Build the raw file manifest}
Create a manifest with one row per discovered file.
For each file, record:
\begin{itemize}
    \item participant folder name;
    \item participant ID extracted from folder;
    \item file name;
    \item file type (\texttt{enem\_blocks}, \texttt{socio\_json}, schema, other);
    \item timestamp extracted from the file name;
    \item full path;
    \item whether the expected pair exists.
\end{itemize}

The purpose of this file is to guarantee that every participant has the expected raw inputs and that
the paired files refer to the same recording timestamp.

\subsubsection{Step~1.2: Build the participant registry}
Create a participant registry with exactly one row per participant.
Since each participant appears only once, the participant registry should include:
\begin{itemize}
    \item \texttt{participant\_id} (e.g. \texttt{PID001});
    \item folder name;
    \item PsychoPy file name;
    \item socio JSON file name;
    \item session timestamp;
    \item matching-status flag;
    \item notes.
\end{itemize}

Even though there is only one session per participant, it is still useful to keep the timestamp,
because later steps will need to connect these files to the corresponding fNIRS recording.

\subsubsection{Step~1.3: Parse the questionnaire schema}
Read \texttt{demographics\_questions.json} and convert it into a flat table called
\texttt{03\_demographics\_schema.csv}.
For each variable, store:
\begin{itemize}
    \item question ID;
    \item question prompt;
    \item declared type;
    \item required flag;
    \item \texttt{show\_if} dependency, if present;
    \item \texttt{show\_if} regular expression, if present.
\end{itemize}

This table becomes the schema reference for all downstream validation and reporting. \cite{step1_demo_schema}

\subsubsection{Step~1.4: Read the raw JSON answers}
Read all \texttt{socio\_PIDxx\_timestamp.json} files into a single raw wide table.
At this stage:
\begin{itemize}
    \item do not recode any value;
    \item do not collapse categories;
    \item do not drop blank strings;
    \item keep exactly what was recorded.
\end{itemize}

This table should be saved as \texttt{04\_socio\_raw\_wide.csv}.

\subsubsection{Step~1.5: Convert the raw wide table to long format}
Create a long-format version with one row per participant-variable pair.
The minimum columns should be:
\begin{verbatim}
participant_id
timestamp
question_id
raw_value
source_json
\end{verbatim}

This long table should be saved as \texttt{05\_socio\_raw\_long.csv}.
The long format will make grouped analysis, quality checks, and automated report generation much easier.

\subsubsection{Step~1.6: Validate against the schema}
For each participant and each schema variable, check:
\begin{itemize}
    \item whether the variable exists in the JSON;
    \item whether the response type is compatible with the schema type;
    \item whether a missing value is allowed;
    \item whether the variable should have been shown according to the \texttt{show\_if} rule.
\end{itemize}

This is very important for conditional variables.
For example:
\begin{itemize}
    \item if \texttt{affiliation} indicates a student, then \texttt{programme\_name},
    \texttt{study\_level}, and \texttt{year\_in\_programme} should be evaluated as relevant fields;
    \item if \texttt{affiliation} indicates an employee, then \texttt{employee\_role}
    becomes relevant;
    \item if a variable was not applicable, it should be coded as \texttt{NOT\_APPLICABLE},
    not as ordinary missing.
\end{itemize}

\subsubsection{Step~1.7: Create the harmonization dictionary}
Create \texttt{06\_socio\_harmonization\_dictionary.csv}.
This file must document every transformation rule.
It should contain at least:
\begin{itemize}
    \item variable name;
    \item raw value;
    \item harmonized value;
    \item transformation rule;
    \item comments;
    \item whether the rule is exact or approximate.
\end{itemize}

No value should be changed without an explicit rule in this dictionary.

\subsubsection{Step~1.8: Apply basic text cleaning}
For all text fields, apply only conservative cleaning:
\begin{itemize}
    \item trim leading and trailing whitespace;
    \item collapse repeated spaces;
    \item standardize obvious capitalization;
    \item preserve accents;
    \item keep the original raw value in a separate column.
\end{itemize}

Examples:
\begin{itemize}
    \item \texttt{" iran "} $\rightarrow$ \texttt{"Iran"}
    \item \texttt{"phd candidate"} $\rightarrow$ \texttt{"PhD candidate"}
    \item \texttt{"udm liser"} $\rightarrow$ \texttt{"UDM LISER"}
\end{itemize}

This step should remain conservative.
No advanced NLP or semantic grouping should be introduced at this stage.

\subsubsection{Step~1.9: Apply variable-specific harmonization}
Each variable should then be harmonized according to its type.

\paragraph{Numeric variables}
The following variables should be parsed as numeric when possible:
\begin{itemize}
    \item \texttt{age}
    \item \texttt{years\_in\_luxembourg}
    \item \texttt{year\_in\_programme}
\end{itemize}

Store:
\begin{itemize}
    \item raw string value;
    \item parsed numeric value;
    \item validity flag.
\end{itemize}

\paragraph{Binary or small-category variables}
The following should be standardized to a small controlled vocabulary:
\begin{itemize}
    \item \texttt{gender}
    \item \texttt{first\_gen\_he}
\end{itemize}

For example:
\begin{itemize}
    \item \texttt{yes}, \texttt{Yes}, \texttt{Y} $\rightarrow$ \texttt{Yes}
    \item \texttt{no}, \texttt{No}, \texttt{N} $\rightarrow$ \texttt{No}
    \item unclear values $\rightarrow$ \texttt{Unknown}
\end{itemize}

\paragraph{Ordered education variables}
The following should be harmonized into ordered levels:
\begin{itemize}
    \item \texttt{study\_level}
    \item \texttt{highest\_degree}
    \item \texttt{parent\_a\_edu}
    \item \texttt{parent\_b\_edu}
\end{itemize}

\paragraph{Open-text fields}
The following should remain essentially open-text, with only light cleaning:
\begin{itemize}
    \item \texttt{country\_of\_birth}
    \item \texttt{nationality}
    \item \texttt{first\_language}
    \item \texttt{other\_languages}
    \item \texttt{department\_unit}
    \item \texttt{programme\_name}
    \item \texttt{degree\_field}
    \item \texttt{country\_secondary\_ed}
    \item \texttt{primary\_school\_langs}
    \item \texttt{living\_arrangement}
\end{itemize}

For these fields, the goal is not to aggressively merge categories at this stage.
The goal is to remove trivial spelling noise while preserving real heterogeneity.

\subsubsection{Step~1.10: Define missingness explicitly}
Use explicit missingness labels:
\begin{itemize}
    \item \texttt{MISSING}: expected but absent;
    \item \texttt{NOT\_APPLICABLE}: hidden by questionnaire logic;
    \item \texttt{INVALID}: present but malformed;
    \item \texttt{UNKNOWN}: ambiguous and unresolved;
    \item \texttt{BLANK}: empty string recorded.
\end{itemize}

This distinction is important for descriptive analysis and later modeling.

\subsubsection{Step~1.11: Build the clean wide table}
Create \texttt{07\_socio\_clean\_wide.csv}, with one row per participant and one harmonized column per variable.
For each variable, it is recommended to keep:
\begin{itemize}
    \item a harmonized analysis column;
    \item the raw value in a parallel column if needed;
    \item a flag column for missingness or uncertainty when useful.
\end{itemize}

\subsubsection{Step~1.12: Produce grouped descriptive analysis for each question}
Each demographic variable should receive its own descriptive analysis block.

For \textbf{numeric variables}, compute:
\begin{itemize}
    \item number of valid responses;
    \item number of missing, invalid, and not-applicable cases;
    \item minimum;
    \item maximum;
    \item mean;
    \item standard deviation;
    \item median;
    \item quartiles.
\end{itemize}

For \textbf{categorical variables}, compute:
\begin{itemize}
    \item number of valid responses;
    \item number of missing, invalid, and not-applicable cases;
    \item frequency of each harmonized category;
    \item percentage of each harmonized category;
    \item number of unique raw forms before harmonization.
\end{itemize}

For \textbf{free-text variables}, compute:
\begin{itemize}
    \item number of valid responses;
    \item number of unique harmonized responses;
    \item top categories by frequency;
    \item optional grouped table if a simple manual recoding exists.
\end{itemize}

\subsubsection{Step~1.13: Export per-question result tables}
The per-question analysis should be saved into at least three files:

\paragraph{\texttt{08\_question\_summary.csv}}
One row per variable, containing:
\begin{itemize}
    \item question ID;
    \item question prompt;
    \item variable group;
    \item variable type;
    \item number of participants;
    \item number of valid responses;
    \item number missing;
    \item number not applicable;
    \item number invalid;
    \item number of unique harmonized categories;
    \item notes.
\end{itemize}

\paragraph{\texttt{09\_question\_level\_counts.csv}}
One row per variable-category pair, containing:
\begin{itemize}
    \item question ID;
    \item harmonized category;
    \item count;
    \item percentage.
\end{itemize}

\paragraph{\texttt{10\_question\_missingness.csv}}
One row per variable and missingness type, containing:
\begin{itemize}
    \item question ID;
    \item missingness label;
    \item count;
    \item percentage.
\end{itemize}

\subsubsection{Step~1.14: Produce plots}
Each variable should have its own plot.

\paragraph{Numeric variables}
Generate:
\begin{itemize}
    \item histogram;
    \item boxplot if useful.
\end{itemize}

\paragraph{Categorical variables}
Generate:
\begin{itemize}
    \item bar plot of harmonized category counts;
    \item optionally a second bar plot for missingness profile.
\end{itemize}

\paragraph{Open-text variables}
Generate:
\begin{itemize}
    \item bar plot of the most frequent cleaned responses;
    \item truncate to top categories when necessary.
\end{itemize}

Plots should be saved in:
\begin{verbatim}
figures/step1/
\end{verbatim}

\subsubsection{Step~1.15: Generate one \LaTeX{} file per variable}
For each variable, generate a separate \LaTeX{} file:
\begin{verbatim}
reports/step1/questions/question_<id>.tex
\end{verbatim}

Each file should contain:
\begin{itemize}
    \item subsection title with question prompt;
    \item variable metadata;
    \item short text summary;
    \item table of counts or summary statistics;
    \item inserted plot;
    \item notes on harmonization and missingness.
\end{itemize}

This modular structure allows the report to be rebuilt automatically without rewriting text by hand.

\subsubsection{Step~1.16: Generate a master \LaTeX{} report}
Create a master file:
\begin{verbatim}
reports/step1/step1_demographics_report.tex
\end{verbatim}

This file should:
\begin{itemize}
    \item introduce the demographic module;
    \item describe the data sources;
    \item summarize the harmonization process;
    \item include grouped sections;
    \item load each variable file with \texttt{\textbackslash input\{...\}}.
\end{itemize}

This makes the report easy to update whenever the data are refreshed.

\subsection{Granular Reporting Template for One Variable}
Each variable-specific \LaTeX{} file should follow the same logic:

\begin{verbatim}
\subsection{Age}
\label{sec:q_age}

\paragraph{Question prompt.}
Age (in years).

\paragraph{Type.}
Numeric, required.

\paragraph{Availability.}
N participants, N valid, N missing, N invalid.

\paragraph{Summary statistics.}
Mean, SD, median, min, max, quartiles.

\paragraph{Quality notes.}
Any implausible values, formatting corrections, or exclusions.

\begin{figure}[ht]
    \centering
    \includegraphics[width=0.7\textwidth]{figures/step1/age_histogram.png}
    \caption{Distribution of age.}
\end{figure}
\end{verbatim}

For a categorical variable, the structure should be analogous, replacing summary statistics with a frequency table.

\subsection{Quality-Control Checks}
Before considering Step~1 complete, the following checks should be run:

\begin{enumerate}
    \item every participant folder has both expected files;
    \item file timestamps match within participant;
    \item schema variables are recognized;
    \item no unexpected variables remain unmapped;
    \item numeric fields parse correctly where expected;
    \item conditional variables are correctly marked as applicable or not applicable;
    \item frequencies are reviewed for suspicious spelling variants;
    \item the final participant count matches the raw folder count.
\end{enumerate}

\subsection{Special Note on Legacy Sessions}
The current questionnaire schema is text-based and uses named variables directly.
However, earlier versions may have stored some answers as numeric option codes rather than text labels.
If such legacy sessions exist, they must not be harmonized by guesswork.
Instead, they should be decoded only from the original option mapping used at the time of collection.
Until that mapping is available, such responses should be flagged as unresolved rather than forced into categories.

\subsection{Minimal Completion Criteria for Step~1}
Step~1 is complete when:
\begin{itemize}
    \item all participant folders have been inventoried;
    \item all socio JSON files have been loaded;
    \item all variables have been checked against the questionnaire schema;
    \item a clean participant-level demographic table exists;
    \item grouped descriptive CSV outputs exist;
    \item one plot per variable exists;
    \item one \LaTeX{} file per variable exists;
    \item one master demographic report exists;
    \item all harmonization decisions are documented.
\end{itemize}

\subsection{Practical Interpretation}
In simple terms, Step~1 should deliver one well-organized demographic analysis package:
a clean participant table, a schema table, harmonization logs, grouped summaries, plots, and a modular
\LaTeX{} report that can be updated automatically whenever new participants are added.

\end{document}
